\documentclass[../main.tex]{subfiles}

\begin{document}

\chapter*{Glossary}
\addcontentsline{toc}{chapter}{Glossary}
\label{app:glossary}

This glossary defines key terms used throughout this book on optimization for machine learning. For foundational mathematical concepts (vectors, matrices, eigenvalues, etc.), please refer to the glossary in \emph{Book 1: The Math That Powers AI}.

\bigskip

\begin{description}[style=nextline, leftmargin=2cm, labelwidth=1.8cm]

\item[Adam]
Adaptive Moment Estimation. An optimizer that combines momentum with adaptive learning rates by maintaining exponential moving averages of both gradients (first moment) and squared gradients (second moment). Update rule: $m_t = \beta_1 m_{t-1} + (1-\beta_1)g_t$, $v_t = \beta_2 v_{t-1} + (1-\beta_2)g_t^2$, with bias correction.

\item[AdaGrad]
Adaptive Gradient Algorithm. An optimizer that adapts learning rates based on cumulative squared gradients, giving smaller updates to frequently occurring features. Effective for sparse data but learning rates can decay too quickly.

\item[Batch Normalization]
A technique that normalizes layer inputs across a mini-batch during training, reducing internal covariate shift and allowing higher learning rates.

\item[Batch Size]
The number of training samples used to compute one gradient estimate. Larger batches reduce variance but increase computational cost per update.

\item[Condition Number]
For a matrix $A$, the ratio $\kappa = \lambda_{\max}/\lambda_{\min}$ of largest to smallest eigenvalue. For optimization, it measures how much the curvature varies across directions; high condition numbers lead to slow convergence.

\item[Convergence Rate]
The speed at which an optimization algorithm approaches the optimum. Linear (geometric) convergence: $\|x_k - x^*\| \leq C\rho^k$ for $\rho < 1$. Sublinear convergence: $f(x_k) - f^* \leq C/k^\alpha$.

\item[Convex Function]
A function $f$ satisfying $f(\alpha x + (1-\alpha)y) \leq \alpha f(x) + (1-\alpha)f(y)$ for all $x, y$ and $\alpha \in [0,1]$. Equivalently, $f$ lies below all its secant lines.

\item[Convex Set]
A set $C$ where for any $x, y \in C$ and $\alpha \in [0,1]$, the point $\alpha x + (1-\alpha)y \in C$. Geometrically, the line segment between any two points lies entirely in the set.

\item[Cross-Entropy Loss]
A loss function derived from maximum likelihood for classification: $\mathcal{L} = -\sum_k y_k \log p_k$ where $y$ is the true label (one-hot) and $p$ is the predicted probability distribution.

\item[Descent Direction]
A direction $d$ along which the function decreases: $\nabla f(x)^\top d < 0$. Gradient descent uses $d = -\nabla f(x)$.

\item[Descent Lemma]
For $L$-smooth functions: $f(y) \leq f(x) + \nabla f(x)^\top(y-x) + \frac{L}{2}\|y-x\|^2$. This quadratic upper bound is fundamental to convergence analysis.

\item[Dropout]
A regularization technique that randomly sets activations to zero during training with probability $p$. At test time, activations are scaled by $(1-p)$ or training uses inverted dropout with scaling $1/(1-p)$.

\item[Early Stopping]
A regularization technique that stops training when validation performance stops improving. Implicitly limits model complexity by constraining the optimization trajectory.

\item[Epoch]
One complete pass through the entire training dataset. Training typically requires multiple epochs.

\item[First-Order Method]
An optimization algorithm using only gradient information (no Hessian). Examples: gradient descent, SGD, Adam.

\item[Gradient]
The vector of partial derivatives $\nabla f(x) = (\partial f/\partial x_1, \ldots, \partial f/\partial x_n)^\top$. Points in the direction of steepest ascent.

\item[Gradient Clipping]
A technique to prevent exploding gradients by scaling gradients when their norm exceeds a threshold: $g \leftarrow g \cdot \min(1, \tau/\|g\|)$.

\item[Gradient Descent]
An iterative optimization algorithm: $x_{k+1} = x_k - \eta \nabla f(x_k)$. Moves in the direction of steepest descent with step size $\eta$.

\item[Hessian]
The matrix of second partial derivatives: $H_{ij} = \partial^2 f/\partial x_i \partial x_j$. Captures local curvature information.

\item[L1 Regularization]
Also called Lasso. Adds penalty $\lambda\|w\|_1$ to the loss. Promotes sparsity by driving some weights exactly to zero.

\item[L2 Regularization]
Also called Ridge regression. Adds penalty $\lambda\|w\|_2^2$ to the loss. Shrinks weights toward zero without sparsity. Note: Often conflated with weight decay, but they differ for adaptive optimizers like Adam---see AdamW.

\item[Learning Rate]
The step size $\eta$ in gradient-based optimization. Controls how much parameters change per update. Also called step size.

\item[Learning Rate Schedule]
A function $\eta(t)$ that varies the learning rate over training. Common schedules: step decay, exponential decay, cosine annealing, warmup.

\item[Learning Rate Warmup]
Starting with a small learning rate and gradually increasing it over initial iterations. Helps stabilize training, especially with large batch sizes or adaptive optimizers.

\item[Lipschitz Continuous]
A function $f$ is $L$-Lipschitz if $|f(x) - f(y)| \leq L\|x - y\|$ for all $x, y$. For gradients: $\|\nabla f(x) - \nabla f(y)\| \leq L\|x - y\|$ (also called $L$-smooth).

\item[Local Minimum]
A point $x^*$ where $f(x^*) \leq f(x)$ for all $x$ in some neighborhood. For non-convex functions, gradient methods may converge to local rather than global minima.

\item[Loss Function]
A function measuring the discrepancy between predictions and targets. Also called objective function or cost function.

\item[Mini-batch]
A subset of training data used to compute a stochastic gradient estimate. Typical sizes: 32, 64, 128, 256.

\item[Momentum]
A technique that accumulates past gradients to accelerate convergence: $v_{t+1} = \beta v_t + \nabla f(x_t)$, $x_{t+1} = x_t - \eta v_t$. Helps navigate ravines and overcome local noise.

\item[MSE Loss]
Mean Squared Error: $\mathcal{L} = \frac{1}{n}\sum_i (y_i - \hat{y}_i)^2$. Standard loss for regression.

\item[Nesterov Momentum]
An accelerated variant of momentum that evaluates the gradient at a ``lookahead'' position: $x_{t+1} = x_t + \beta(x_t - x_{t-1}) - \eta\nabla f(x_t + \beta(x_t - x_{t-1}))$. Achieves optimal convergence rates for convex functions.

\item[Optimizer]
An algorithm for minimizing a loss function by updating model parameters. Examples: SGD, Adam, RMSprop.

\item[Overfitting]
When a model performs well on training data but poorly on unseen data. Indicates the model has memorized training examples rather than learning generalizable patterns.

\item[Proximal Operator]
The operator $\text{prox}_{\lambda g}(x) = \arg\min_z \{g(z) + \frac{1}{2\lambda}\|z-x\|^2\}$. Used in proximal gradient methods for non-smooth optimization.

\item[Regularization]
Techniques to prevent overfitting by constraining model complexity. Includes: L1/L2 penalties, dropout, early stopping, data augmentation.

\item[RMSprop]
Root Mean Square Propagation. An optimizer that adapts learning rates using an exponential moving average of squared gradients: $v_t = \beta v_{t-1} + (1-\beta)g_t^2$, $x_{t+1} = x_t - \frac{\eta}{\sqrt{v_t + \epsilon}}g_t$.

\item[Saddle Point]
A critical point where $\nabla f(x) = 0$ but the Hessian has both positive and negative eigenvalues. Common in high-dimensional non-convex optimization.

\item[Second-Order Method]
An optimization algorithm using Hessian information. Examples: Newton's method, quasi-Newton methods (L-BFGS). Faster convergence but higher per-iteration cost.

\item[SGD]
Stochastic Gradient Descent. Approximates the full gradient using a random subset of data: $x_{k+1} = x_k - \eta \nabla f_i(x_k)$ where $i$ is randomly sampled.

\item[Smoothness]
A function is $L$-smooth if its gradient is $L$-Lipschitz continuous. Equivalent to: $f(y) \leq f(x) + \nabla f(x)^\top(y-x) + \frac{L}{2}\|y-x\|^2$.

\item[Stationary Point]
A point where $\nabla f(x) = 0$. Includes minima, maxima, and saddle points.

\item[Step Size]
See Learning Rate.

\item[Stochastic Gradient]
A random estimate of the true gradient, typically computed from a mini-batch. Has the same expectation as the true gradient but non-zero variance.

\item[Strong Convexity]
A function $f$ is $\mu$-strongly convex if $f(y) \geq f(x) + \nabla f(x)^\top(y-x) + \frac{\mu}{2}\|y-x\|^2$. Implies a unique minimum and guarantees linear convergence for gradient descent.

\item[Subgradient]
A generalization of gradients to non-differentiable convex functions. Vector $g$ is a subgradient of $f$ at $x$ if $f(y) \geq f(x) + g^\top(y-x)$ for all $y$.

\item[Underfitting]
When a model performs poorly on both training and test data. Indicates insufficient model capacity or inadequate training.

\item[Variance Reduction]
Techniques to reduce the variance of stochastic gradient estimates without increasing batch size. Examples: SVRG, SAGA.

\item[Weight Decay]
A regularization technique that shrinks weights toward zero each update: $w \leftarrow (1-\lambda\eta)w - \eta\nabla\mathcal{L}$. Equivalent to L2 regularization for vanilla SGD but different for adaptive optimizers.

\end{description}

\end{document}
